\section{Анализ предметной области}
\subsection{Файловые системы}

Файловая система -- структура, определяющая порядок организации и хранения данных на носителях информации в ЭВМ. Файловые системы управляют хранением файлов, каталогов и метаданных, определяют способ доступа к данных, их размещение и защиту. Без этого данные представляют собой последовательность байтов без структуры.

Основные функции файловой системы включают:

\begin{itemize}
\item создание файлов;
\item запись в файлы;
\item чтение из файлов;
\item поиск файлов;
\item удаление файлов;
\item отслеживание свободного пространства;
\item предотвращение конфликтов доступа;
\item обеспечение целостности данных.
\end{itemize}

Файловые системы делятся на:

\begin{itemize}
\item дисковые -- на носителях информации(FAT32, CDFS, NTFS, ext4, APFS);
\item сетевые -- сетевой доступ к файловой системе(NFS, SMB);
\item виртуальные -- абстракция поверх конкретной файловой системы, цель которой обеспечить доступ к разным файловым системам.
\end{itemize}

Файловая система включает в себя:

\begin{enumerate}
\item Блок данных -- физический блок информации, в котором хранится файл или его часть. Представляет собой фрагмент данных, который может быть прочитан или записан, чаще всего имеет фиксированный размер.
\item Таблица индексов. Она содержит упорядоченную информацию о расположении блоков данных(индексы), чем ускоряет чтение и поиск файлов.
\item Журналы транзакций -- файл, в котором записаны недавние изменения в файловой системе, такие как создание, удаление или изменение файлов, которые помогают обеспечивать целостность данных и восстановить их в случае сбоев.
\item Метаданные -- информация о атрибутах файла таких как размер файла, его тип, время создания и изменения, права доступа и другие.
\end{enumerate}

\subsection{Типы файловой системы}

FAT(File Allocation Table) представляет собой один из первых типов файловой системы. Microsoft создала ее для floppy-дисков. Число после FAT означает сколько бит выделяется под адрес, например FAT12 использовала 12-битные адреса. FAT подходит для небольших дисков, но имеет ограничения в размере.

NTFS(New Technology File System) пришла на смену FAT в Windows NT. Она поддерживает большие объемы, ACL для безопасности и компрессию. NTFS использует MFT(Master File Table) для хранения метаданных.

В Linux чаще всего используются ext2, ext3(ext2 с журналированием) и ext4(ext3 с улучшенной производительностью). Эти системы используют inode для описания файлов.

Для macOS Apple использует HFS+ и APFS. APFS добавила snapshots и клонирование файлов.

Сетевые файловые системы включают NFS от Sun Microsystems и CIFS/SMB от Microsoft. Они позволяют совместный доступ к файлам через сеть.

Виртуальные файловые системы, как tmpfs, хранят данные в оперативной памяти.

Файловые системы для флэш-памяти, такие как JFFS или YAFFS, учитывают износ ячеек.

Файловые системы для оптических носителей отличаются. CDFS базируется на ISO 9660. Она предназначена для CD-ROM, где данные записывают один раз.

\subsection{История FAT}

FAT появилась в 1977 году с Microsoft Disk BASIC. Первая версия использовала 8-битные адреса.

В 1980 году FAT12 стала стандартом для MS-DOS 1.0. Она поддерживала диски до 16 МБ.

FAT16 вышла в 1984 году с MS-DOS 3.0. Она позволяла использовать диски до 2 ГБ. Кластеры имели размер от 2 КБ до 32 КБ.

FAT32 ввели в 1996 году с Windows 95 OSR2. Она поддерживает диски до 2 ТБ. Кластеры начинаются от 4 КБ. FAT32 добавила резервную копию FAT и улучшила управление пространством на диске.

FAT использовали в USB-накопителях и SD-картах. Ее простота обеспечивает совместимость между операционными системами.

\subsection{Устройство FAT32}

Ключевая концепция FAT32, как и любой FAT, -- это таблица размещения файлов(File Allocation Table). Это карта всего дискового пространства, которая показывает, какие кластеры свободны, какие заняты, а какие являются поврежденными. Каждому файлу и каталогу соответствует цепочка кластеров, связи между которыми хранятся в этой таблице.

\subsubsection{Структура диска с файловой системой FAT32}

Диск с FAT32 делится на 3 основные области:

\begin{enumerate}
\item Зарезервированная область(Reserved Region). Эта область начинается с самого первого сектора диска(0 сектор) и содержит критически важную информацию для системы. Главная составляющая -- загрузочный сектор(BPB -- BIOS Parameter Block). Это первый сектор области. Он содержит служебные данные, необходимые для работы ОС с этой файловой системой. Ключевые поля в BPB:
\begin{enumerate}
\item Размер сектора.
\item Количество секторов в одном кластере. Кластер -- это минимальная единица размещения файла. Размер кластера зависит от размера тома.
\item Количество зарезервированных секторов(включая сам boot-сектор). Для FAT32 обычно равно 32.
\item Количество таблиц FAT. Почти всегда 2(основная и резервная копия для избыточности).
\item Размер одной таблицы FAT в секторах.
\item Номер первого кластера корневого каталога. В FAT32 корневой каталог может быть любого размера и расположен в области данных, а не в фиксированном месте.
\item Сектор в зарезервированной области, где хранится структура FSINFO со служебной информацией(количество свободных кластеров, номер следующего свободного кластера).
\item Сектор, где хранится копия загрузочного сектора.
\end{enumerate}
\item Область FAT(FAT Region). Сразу после зарезервированной области идут подряд несколько копий таблицы FAT(их количество указано в BPB). Сама таблица FAT -- это массив 32-битных элементов(отсюда и название FAT32). Каждому элементу этого массива соответствует один кластер из области данных.
\item Область данных(Data Region). Это основное пространство для хранения файлов и каталогов, разделённое на кластеры, пронумерованные начиная с 2, потому что кластеры 0 и 1 не существуют, их номера используются для служебных целей в FAT. Корневой каталог это не отдельная область, а обычный каталог, который начинается с кластера, номер которого указан в BPB(Root Cluster). Он может расти и располагаться в любом месте области данных, как и любой другой каталог. Подкаталоги представляют собой специальные файлы, содержащие записи о файлах и других каталогах внутри себя.
\end{enumerate}

\subsubsection{Структура записей в каталогах}

Каждый файл и каталог в области данных представлен одной или несколькими 32-байтными записями в родительском каталоге.

Стандартная 32-байтная запись содержит:

\begin{enumerate}
\item Имя файла(8 байт) и его расширение(3 байта). Для длинных имён файлов перед стандартной записью файла размещаются дополнительные записи, связанные в цепочку. Каждая такая LFN-запись может хранить до 13 символов имени в Unicode.
\item Атрибуты(1 байт), в нём содержатся флаги(например, Read-Only, Hidden, System, Volume Label, Directory, Archive).
\item Резервные поля: время создания(2 байта), дата создания(2 байта), дата последнего доступа(2 байта).
\item Время модификации(2 байта)
\item Дата модификации(2 байта)
\item Номер первого кластера(2 байта high 2 bytes). Старшие 2 байта номера первого кластера.
\item Номер первого кластера(2 байта low 2 bytes). Младшие 2 байта номера первого кластера.
\item Размер файла(4 байта). Размер файла в байтах. Для каталогов обычно указан 0.
\end{enumerate}

\subsubsection{Преимущества и недостатки FAT32}

Преимущества файловой системы FAT32:

\begin{itemize}
\item кроссплатформенность(поддерживается практически всеми операционными системами и устройствами);
\item простота(имеет очень простую структуру, что позволяет файловой системе быстро работать на слабых устройствах и встроенных системах);
\item эффективность для носителей с малым объёмом(на флеш-накопителях и картах памяти небольшого объема FAT32 показывает очень хорошую скорость работы);
\item широкая поддержка инструментами(из-за своей распространенности существует огромное количество утилит для восстановления данных с FAT32-носителей).
\end{itemize}

Недостатки файловой системы FAT32:

\begin{itemize}
\item ограничение на максимальный размер файла(4 ГБ);
\item ограничение на максимальный размер тома(раздела);
\item отсутствие журналирования;
\item низкая отказоустойчивость;
\item сильная фрагментация;
\item отсутствие современных функций безопасности и разграничения прав;
\item неэффективное использование пространства на больших дисках.
\end{itemize}

\subsection{Устройство CDFS(ISO 9660)}

CDFS(CD File System) -- является реализацией стандарта ISO 9660 для CD-ROM.

ISO 9660 -- это международный стандарт, разработанный в 1988 году для обеспечения совместимости компакт-дисков между разными операционными системами и аппаратными платформами. Его основная цель — предоставить простую и универсальную файловую систему для CD-ROM.

\subsubsection{Структура диска с файловой системой CDFS}

Диск с CDFS делится на 3 основные области:

\begin{enumerate}
\item Область дескрипторов тома(Volume Descriptors). Эта область находится в самом начале диска(после системной области начиная с 16-го логического сектора) и описывает всё содержимое диска. Она состоит из последовательности дескрипторов, каждый размером ровно в 1 сектор(2048 байт). Дескрипторы действуют как заголовки для диска. Типы дескрипторов:
\begin{enumerate}
\item основной дескриптор тома(Primary Volume Descriptor или PVD). Содержит всю базовую информацию о томе: идентификатор системы, идентификатор тома, размер сектора, количество секторов в томе, расположение корневого каталога(Root Directory Record), расположение таблицы путей(Path Tables) — L-форма и М-форма.
\item ограничитель набора дескрипторов тома(Volume descriptor set terminator). Обозначает конец набора дескрипторов.
\end{enumerate}
\item Таблицы путей(Path Tables). Эта область представляет собой таблицу, содержащую записи каждого каталога на диске. Благодаря этой таблице можно быстро найти местоположение любого каталога. Каждая запись содержит:
\begin{itemize}
\item идентификатор каталога;
\item номер родительского каталога;
\item номер сектора записи каталога в области данных.
\end{itemize}
Таблица путей хранится в 2 экземплярах: L-Path table(порядок байт от младшего к старшему) и M-Path table(порядок байт от старшего к младшему).
\item Область данных. В этой области хранятся данные файлов и записи каталогов.
Каждый файл и подкаталог представлен записью каталогов(Directory Record), который включает:
\begin{itemize}
\item размер записи;
\item расположение экстента(номер сектора, на котором начинаются данные файла или запись каталога);
\item размер файла в байтах;
\item атрибуты(например IsDirectory или Hidden);
\item длина имени файла;
\item имя файла;
\item дата и время записи.
\end{itemize}
В ISO 9660 нет фрагментации, файлы хранятся в виде экстента(последовательной серии секторов), начиная с номера в записи каталога.
\end{enumerate}

\subsubsection{Расширения стандарта ISO 9660}

ISO 9660 имеет множество ограничений, из-за чего были созданы расширения, сосуществующие с основной структурой:

\begin{enumerate}
\item Joliet. Создан Microsoft для поддержки длинных имён файлов в Unicode. Реализовано через создание нового типа дескриптора тома -- Supplementary Volume Descriptor,
который дополнительно указывает на свою копию корневого каталога и таблиц путей с именами в Unicode. Системы, не поддерживающие Joliet игнорируют SVM,
а поскольку Joliet сохраняет файловую систему ISO 9660, никаких проблем с совместимостью не возникает.
\item Rock Ridge Interchange Protocol(RRIP). Создан для поддержки файловых систем UNIX(POSIX): длинные имена файлов, права доступа, идентификаторы пользователя и группы.
Реализован через добавления дополнительных полей в запись каталога, которые игнорируется системами без этого расширения.
\item El Torito. Цель -- создание загрузочных CD-ROM. Реализовано через создание нового типа дескриптора -- Boot Record Descriptor, который содержит информацию,
эмулирующую жёсткий диск или флоппи-диск, с которых BIOS может загрузить информацию загрузочного диска.
\end{enumerate}

\subsubsection{Преимущества и недостатки CDFS}

Преимущества файловой системы CDFS:

\begin{itemize}
\item кроссплатформенность(поддерживается практически всеми операционными системами и устройствами);
\item простота(для этой файловой системы легко написать драйвер);
\item эффективность для последовательного доступа;
\item таблицы путей позволяют быстро находить каталоги;
\item поддержка расширений, дополняющих ISO 9660.
\end{itemize}

Недостатки файловой системы CDFS:

\begin{itemize}
\item без расширений имеет устаревшие и строгие ограничения;
\item только для чтения;
\item несмотря на таблицы путей, из-за особенностей CD-приводов произвольный доступ к мелким файлам может быть очень медленным;
\item избыточность данных;
\item нет шифрования.
\end{itemize}

\subsection{История ext(Extended File System)}

Файловая система ext появилась в 1992 году, и является первой файловой системой, специально созданной для ядра Linux.

ext2 вышла в 1993 для решения множества проблем ext.

ext3, вышедшая в 2001 году, добавила в ext2 поддержку журналирования, которое имело 3 режима:

\begin{itemize}
\item writeback -- сохраняются только метаданные;
\item ordered -- writeback, но информация об изменении файла происходит после записи в файл;
\item journal -- сохраняются метаданные файлов и их содержимое.
\end{itemize}

ext4, созданный в 2008 году, в отличии от ext3:

\begin{itemize}
\item увеличил размер одного раздела диска;
\item увеличил максимальный размер файла;
\item добавил возможность объединять блоки в экстенты при записи файла.
\end{itemize}

\subsection{Устройство ext2}

Ключевая концепция ext2(Second Extended File System) -- разделение пространства диска на блоки, которые объединяются в группы блоков.
Каждая группа блоков хранит информацию про занятость блоков внутри группы, а также метаданные файлов и их содержимое.

\subsubsection{Структура диска с файловой системой ext2}

Первый сектор диска с файловой системой ext2 является загрузчиком операционной системы. Далее диск разделён на блоки, а блоки объединены в группы блоков.
Каждая группа блоков состоит из нескольких частей:

\begin{enumerate}
\item Суперблок(обязательно находится в 0 группе блоков, в следующих группах блоков могут храниться копии). Размер суперблока 1 КБ. Он содержит:
\begin{itemize}
\item магическое число -- идентификатор файловой системы;
\item количество блоков и индексных дескрипторов(inode) в файловой системе;
\item количество свободных блоков и индексных дескрипторов в файловой системе;
\item размер одного блока;
\item размер одного индексного дескриптора;
\item количество блоков в группе;
\item номер первого блока(принимает значение 0, если размер блока больше 1 КБ, и 1, если равен 1 КБ).
\end{itemize}
\item Дескриптор группы блоков(обязательно находится в 0 группе блоков, в следующих группах блоков могут храниться копии) -- массив, в котором каждый элемент описывает одну группу блоков.
Каждый дескриптор группы блоков содержит:
\begin{itemize}
\item номер блока, где находится битовая карта занятости блоков группы;
\item номер блока, где находится битовая карта занятости индексных дескрипторов этой группы;
\item номер блока, где начинается таблица индексных дескрипторов группы;
\item количество свободных блоков в группе;
\item количество свободных индексных дескрипторов в группе;
\item количество каталогов в группе.
\end{itemize}
\item Битовая карта занятости блоков(обязательно во всех группах блоков). Каждый бит -- блок данных в группе. Принимает значение 1, если блок занят, и 0, если свободен.
\item Битовая карта занятости индексных дескрипторов(обязательно во всех группах блоков). Каждый бит -- индексный дескриптор в группе. Принимает значение 1, если индексный дескриптор занят, и 0, если свободен.
\item Таблица индексных дескрипторов(обязательно во всех группах блоков) -- массив, в котором каждый элемент описывает один файл или каталог.
Каждый индексный дескриптор содержит:
\begin{itemize}
\item тип файла;
\item права доступа;
\item идентификатор владельца;
\item размер файла;
\item время последнего доступа;
\item время последней модификации;
\item время последнего изменения;
\item количество выделенных блоков;
\item массив из 15 указателей на блоки, где хранятся данные файла или каталога.
\end{itemize}
\item Блоки данных(обязательно во всех группах блоков) -- последовательность блоков, в которых хранятся данные файлов и каталогов.
\end{enumerate}

\subsubsection{Организация данных файла в ext2}

В файловой системе ext2, в каждом индексном дескрипторе хранится массив из 15 элементов -- многоуровневая схема для эффективного хранения как маленьких, так и больших файлов.
Массив содержит:
\begin{enumerate}
\item Прямая ссылка. Занимает 12 элементов массива и содержит номера блоков, в которых хранятся данные файла.
\item Косвенная ссылка. Занимает 13 элемент массива и содержит номер блока, который является массивом номеров блоков, в которых хранятся данные файла.
\item Двойная косвенная ссылка. Занимает 14 элемент массива и содержит номер блока, который является массивом косвенных ссылок.
\item Тройная косвенная ссылка. Занимает 15 элемент массива и содержит номер блока, который является массивом двойных косвенных ссылок.
\end{enumerate}

\subsubsection{Каталог в ext2}

Каталог в ext2 -- это файл, содержимое которого представляет собой список записей каталога.
Каждая запись каталога содержит:
\begin{itemize}
\item длину этой записи каталога;
\item номер индексного дескриптора файла;
\item длину имени файла;
\item массив переменной длины заканчивающийся 0, в котором хранится имя файла.
\end{itemize}

\subsubsection{Преимущества и недостатки ext2}

Преимущества файловой системы ext2:

\begin{itemize}
\item высокая производительность;
\item поддерживается большинством Linux-дистрибутивов;
\item контроль доступа(ACL) и расширенных атрибутов файлов по стандарту POSIX.
\item эффективность для устройств хранения данных на основе флэш-памяти.
\end{itemize}

Недостатки файловой системы ext2:

\begin{itemize}
\item ограничение размера;
\item отсутствие журналирования;
\item обязательная проверка(после некорректного завершения работы) на дисках большого объёма может занимать много времени;
\item нет шифрования.
\end{itemize}